\chapter{Glossary of school terms}
\label{ch:terminologia_didattica}
\index{Glossary!of school terms}
\index{Terminology!didactic}

\epigraph{``Words are things, and a small drop of ink, falling
like dew upon a thought, produces that which makes thousands,
perhaps millions, think.''}{\textsc{Lord Byron}, \emph{Don Juan},
III, 88}

\lettrine[lines=3]{T}{he} timetabling of Italian schools carries
with it a rich lexicon that is often not shared with the rest of
the world. One and the same entity is called \emph{cattedra}
(teaching post) by the principal, \emph{assegnamento} (assignment)
by the information system, \emph{ore di servizio} (service hours)
by the union, \emph{tuples $(t,c,s)$} by the solver. This chapter
fixes the vocabulary: each entry carries the canonical term, the
aliases, an operational definition and a pointer to the database
tables or the code modules where the concept lives.

\section{Registry entities}

\begin{description}
\item[Docente (Teacher).]\index{Teacher}\index{Cattedra}
  \emph{Aliases}: insegnante, prof, teacher (code). A physical
  person holding one or more \emph{cattedre} (teaching posts)
  within the school. In the database it is a row of the
  \texttt{teachers} table, with \texttt{name}, \texttt{group} (the
  \emph{classe di concorso}, the competition class / subject
  qualification), weekly \texttt{max\_hours} and a set of flags
  (sostegno, potenziamento). Its unavailabilities are in
  \texttt{teacher\_unavailability}.

\item[Cattedra (Teaching post).]\index{Cattedra}\index{Assignment}
  \emph{Aliases}: assegnamento (assignment), cattedra di
  insegnamento, row of the \texttt{Assignment}. A triple
  $(\text{teacher},\text{class},\text{subject})$ with the number
  of weekly hours. A mathematics cattedra for 1A worth 5 hours is
  the row $(\text{Borghi}, 1A, \text{Matematica}, 5)$ in the
  \texttt{assignments} table. All the hours of a cattedra end up
  in the six days of the week and in the six time slots of each
  day: the constraint $\sum_{d,h} x_{tcsdh} = \text{hours}$ is the
  heart of the CP-SAT model.

\item[Classe (Class).]\index{Class}
  \emph{Aliases}: teaching room (sometimes), ``the class 1A''.
  It is a group of students who share the study programme. In the
  database it is a row of \texttt{classes} with \texttt{name}
  (e.g.\ \texttt{1A}), \texttt{year} (1\ldots5), \texttt{section}
  (A, B, C\ldots), \texttt{curriculum} (e.g.\ Scientific,
  Linguistic) and \texttt{n\_students}. The distinction between
  \emph{class} and \emph{classroom} is essential: the class is a
  set of students, the classroom is a physical unit.

\item[Aula (Classroom).]\index{Classroom}
  \emph{Aliases}: laboratory, gym, classroom (code). A physical
  space with characteristics (capacity, equipment, plesso it
  belongs to). Table \texttt{classrooms}. The constraint \emph{the
  lessons of a class in one slot must take place in the same
  classroom} is an implicit standard (with exceptions for
  articulations and laboratories).

\item[Materia (Subject).]\index{Subject}
  \emph{Aliases}: subject, discipline. The subject in itself is
  simply a name (\emph{Matematica} / Mathematics, \emph{Italiano}
  / Italian). The hours of each subject in each class are the
  Cartesian product (class $\times$ subject $\to$ hours) that
  lives in \texttt{class\_subjects}.

\item[Studente (Student).]\index{Student}
  Table \texttt{students}. For the base model $\pi$Tantum does not
  place individual students into slots --- it does so for classes.
  But when there are \emph{articulated groups} or
  \emph{potenziamenti} (enrichment courses), the student becomes a
  first-class entity because their \texttt{tags} and
  \texttt{groups} influence the aggregation.
\end{description}

\section{Organizational structures}

\begin{description}
\item[Plesso (Campus).]\index{Plesso}
  \emph{Aliases}: succursale, detached site. A separate physical
  building under the same school. A school with three plessi has
  three logical coordinates $(x,y)$ in \texttt{plessi}. A typical
  constraint: ``Borghi cannot have one hour at Plesso A at 9 and
  one hour at Plesso B at 10'' (impossible commuting). It is
  handled by the \texttt{plesso\_commuting\_rules} table with a
  minimum \texttt{min\_gap\_hours} between two hours in different
  plessi.

\item[Indirizzo (Study track / curriculum).]\index{Indirizzo}\index{Curriculum}
  \emph{Aliases}: curriculum, profile, thematic section. It
  defines the \emph{study plan}: which subjects and how many
  hours. A \emph{liceo scientifico} (science high school) has
  Latin, a \emph{liceo scienze applicate} (applied-sciences high
  school) does not; a \emph{scientifico} does more hours of
  mathematics than a \emph{linguistico} (language high school).
  Table \texttt{curricula}. The data model allows curricula to
  share subjects, and a single curriculum to change the hours per
  subject per year (1A scientifico does 5 hours of mathematics, 3A
  scientifico does 4).

\item[Anno (Year).]\index{School year}
  \emph{Aliases}: year class, year of course. In Italian licei and
  institutes it is an integer in $\{1,2,3,4,5\}$. The classes of
  the \emph{biennio} (first two years: 1, 2) tend to be more
  numerous (5--6 sections); those of the \emph{triennio} (last
  three years: 3, 4, 5) sparser. In code the field is
  \texttt{class.year}.

\item[Sezione (Section).]\index{Section}
  \emph{Aliases}: corso (stream). The letter of the class (1A, 1B,
  1C\ldots). It is the \emph{third} index (year, section,
  curriculum) that uniquely identifies a class.
\end{description}

\section{Special teaching forms}

\begin{description}
\item[Compresenza (Co-teaching).]\index{Compresenza}\index{Coteach}
  \emph{Aliases}: co-teaching, co-docenza, coteach (code). Two
  teachers simultaneously on the same class in the same hour,
  usually for subjects with a strong laboratory component
  (language $+$ native-speaker conversation assistant; science $+$
  laboratory technician). Table \texttt{coteach\_groups} with
  \texttt{required=True} for the mandatory hours,
  \texttt{required=False} for the preferred hours. The number of
  co-teaching hours is \texttt{n\_hours}; the CP-SAT constraint is
  $x_{t_1} = x_{t_2}$ in all the designated slots.

\item[Sostegno (Special-needs support).]\index{Sostegno}
  \emph{Aliases}: docente di sostegno, BES support. A specialized
  teacher who follows one or more students with special
  educational needs, usually in co-teaching with the subject
  teacher. Their presence is modelled as \texttt{is\_support=True}
  on the corresponding \texttt{assignments} row, and the slots in
  which they are co-present coincide with those of the other
  teacher.

\item[Potenziamento (Enrichment hours).]\index{Potenziamento}
  \emph{Aliases}: hours of empowerment, ore di potenziamento.
  Additional hours of a subject beyond the ministerial minimum,
  usually delivered by a second teacher of the same competition
  class. It is \texttt{is\_potenziamento=True} on
  \texttt{assignments} and typically does \emph{not} translate
  into co-teaching hours: they are separate hours.

\item[Gruppo articolato (Articulated group).]\index{Articulated group}
  \emph{Aliases}: classe articolata, articulated class, group
  (code). A subset of a class with a specific tag (for example:
  the students of second language French vs.\ Spanish within 3A
  linguistico). For the solver a group is a separate entity that
  lives in the same slot as its base class. Table \texttt{groups};
  the articulated cattedre have a non-null \texttt{group\_id}.
  Typical constraint: ``in 1A linguistico, the hour of French
  conversation and the hour of Spanish conversation go in the same
  slot but in two different classrooms''.

\item[Parallel group.]\index{Parallel group}
  Cattedre of different classes fixed at the same hour, usually
  because the students are split across transversal groups (e.g.\
  physical education divided into male/female from two parallel
  classes). Table \texttt{parallel\_groups}; constraint
  $x_{t_1, c_1, s, d, h} = x_{t_2, c_2, s, d, h}$ for each
  $(d, h)$ of the group.
\end{description}

\section{Time slots and calendar}

\begin{description}
\item[Slot.]\index{Slot}
  A pair $(\text{day}, \text{hour})$ in
  $\{1,\ldots,6\} \times \{1,\ldots,6\}$ for the base model. That
  is 36 slots per week. A lesson occupies exactly one slot.

\item[Giorno (Day).]
  Index from 1 (Monday) to 6 (Saturday) in the classic Italian
  model. Saturday is configurable as \emph{free} for schools that
  adopt the short week.

\item[Ora (Hour).]
  Index from 1 to 6, corresponding typically to 8:00, 9:00,
  10:00, 11:00, 12:00, 13:00. The maximum number of hours per day
  is configurable per profile and per track.

\item[Sesta ora (Sixth hour).]\index{Sixth hour}
  The 6\textsuperscript{th} hour (\texttt{hour=13} in the legacy
  code, \texttt{6} in the current model) is penalized because in
  many institutes it ends at 13:50 and is preferably avoided.

\item[Buco (Gap).]\index{Gap}
  An empty slot between two full slots in the same day of the same
  teacher or the same class. It is penalized in the soft score.

\item[Working hours.]\index{Working hours}
  Table \texttt{working\_hours} (imported from the pickles): for
  each (day, hour) a flag that says whether that slot is
  \emph{operational} or \emph{closed}. Closed slots receive no
  lessons. It allows modelling schools with a reduced Saturday
  timetable (5 hours) or without Saturday.

\item[Holiday calendar.]\index{Calendar}\index{Holidays}
  Overlaid on the model week, it describes the dates on which the
  school is closed (holidays, bridge days, strikes). It does not
  enter the CP-SAT model directly: it serves only for the XLSX
  export and the Schedule view to hide the non-working days.
\end{description}

\section{Constraints and preferences}

\begin{description}
\item[Hard.]\index{Constraint!hard}
  A constraint that the solution \emph{must} respect. A hard
  violation makes the solution infeasible: the Feasibility Check
  marks it as such and the pipeline fails or falls back.

\item[Soft.]\index{Constraint!soft}
  A constraint that one \emph{would prefer} to respect but that
  can be violated by paying a penalty (the
  \texttt{soft\_penalty}). The softs feed the objective that
  CP-SAT minimizes.

\item[Preferred.]\index{Constraint!preferred}
  Equivalent to soft with a \emph{very low} penalty: the solver
  satisfies it only if free.

\item[Enforced.]\index{Constraint!enforced}
  Equivalent to hard with an explicit enforcement priority: it
  represents ``institutional'' constraints (usually of legal
  origin).

\item[Indisponibilit\`a (Unavailability).]\index{Unavailability}
  A slot in which an entity (teacher, class, classroom) cannot be
  used. The dedicated tables are
  \texttt{teacher\_unavailability},
  \texttt{class\_unavailability},
  \texttt{classroom\_unavailability}, each with \texttt{state}
  $\in$ \{hard, soft\}.

\item[Logical unavailability.]\index{Logical unavailability}
  A synthetic form for complex unavailabilities:
  \texttt{logical\_unavailabilities} admits a DNF expression
  (e.g.\ ``the teacher is not available on Monday morning OR
  Tuesday afternoon''). It is derived automatically from the DSL.

\item[5 states.]\index{Five states}
  The tables \texttt{teacher\_class\_preferences} and
  \texttt{teacher\_curriculum\_preferences} use a 5-level scale:
  \texttt{forbidden} (hard excluded), \texttt{soft} (medium
  penalty), \texttt{allowed} (default), \texttt{preferred}
  (bonus), \texttt{enforced} (hard imposed). It is the
  formalization of the classic preferences of a school principal.
\end{description}

\section{Terms of the mathematical model}

\begin{description}
\item[$x_{tcsdh}$.]\index{Boolean variable}
  Binary variable of the CP-SAT. It is 1 if in slot $(d,h)$ the
  teacher $t$ teaches the subject $s$ to the class $c$. It is the
  \emph{heart} of the model.

\item[$\text{dc}_{tcsd}$.]\index{Day count}\index{dc value}
  Daily count of hours: it lies between 0 and
  $\text{hours}_{tcs}$. Produced by Phase A; consumed as ``floor''
  or ``exact value'' by Phase B. It is not a classic decision
  variable of the final model, but a decision already taken in the
  previous phase.

\item[Master LP.]\index{Master LP}
  In the Column Generation framework it is the relaxed problem
  that decides \emph{which patterns to select} with non-integer
  coefficients. The dual variables of the master feed the pricers.

\item[Pattern (column).]\index{Pattern}\index{Column}
  A feasible solution restricted to one entity (usually: a
  teacher). The pricer generates it, the master evaluates it, and
  the final $\lambda$-combination is the timetable. A column in
  branch-and-price is encoded as a set of slots
  $\{(c,s,d,h)\}$.

\item[Reduced cost.]\index{Reduced cost}
  For a candidate column, reduced cost $=$ cost of the pattern $-$
  contribution of the duals. A column enters the master if it has
  reduced cost $<\,$0.

\item[Ryan-Foster branching.]\index{Ryan-Foster}
  A branching strategy on pairs of patterns: given a fractional
  $\lambda^* \in (0,1)$, one looks for a pair of slots that some
  patterns use together and others do not, and branches by
  imposing or forbidding their simultaneous use.

\item[Dual stabilization.]\index{Dual stabilization}
  A technique to avoid oscillations of the duals between
  iterations of Column Generation: the current dual vector is
  mixed with the previous one, smoothing the variations.
\end{description}

\section{Pricing granularity in branch-and-price}

\index{Granularity!pricing}

The \texttt{granularity=} parameter of
\texttt{run\_column\_generation()}
(\texttt{engine/column\_generation.py}) admits nine values ---
from \texttt{teacher} (the coarsest, one column per teacher,
default) up to \texttt{day} and \texttt{curriculum} (the finest,
one column per day or per track) --- each with its own enumeration
function \texttt{\_enumerate\_pricing\_keys()}. The complete
taxonomy, with the columnar encoding and the indication of when to
prefer one granularity over another, is in
Table~\ref{tab:bp-granularity-en} of
chapter~\ref{ch:tecniche_avanzate_en}.

\fleuron

To these entries are added the algorithmic terms
(LNS, ALNS, SA, TS, ILS, VNS, Hall, Lagrangian, Column
Generation, Monte Carlo) treated in
chapter~\ref{ch:glossario_discorsivo_en} in a narrative language.
The present chapter does not repeat them: as a rule, if a term is
not found here, it should be looked for there.
